\documentclass[12pt, a4paper]{article}

% ── Packages ─────────────────────────────────────────────────
\usepackage[utf8]{inputenc}
\usepackage[T1]{fontenc}
\usepackage{amsmath, amssymb, amsthm}
\usepackage{mathtools}
\usepackage{graphicx}
\usepackage[margin=1in]{geometry}
\usepackage{hyperref}
\usepackage{cleveref}
\usepackage{booktabs}
\usepackage{array}
\usepackage{caption}
\usepackage{subcaption}
\usepackage{enumitem}
\usepackage{microtype}
\usepackage{xcolor}

\hypersetup{
  colorlinks=true,
  linkcolor=blue!60!black,
  citecolor=green!40!black,
  urlcolor=blue!60!black,
}

% ── Theorem environments ─────────────────────────────────────
\newtheorem{theorem}{Theorem}[section]
\newtheorem{proposition}[theorem]{Proposition}
\newtheorem{lemma}[theorem]{Lemma}
\newtheorem{corollary}[theorem]{Corollary}
\theoremstyle{definition}
\newtheorem{definition}[theorem]{Definition}
\theoremstyle{remark}
\newtheorem{remark}[theorem]{Remark}

% ── Custom commands ──────────────────────────────────────────
\newcommand{\Sthree}{S^3}
\newcommand{\Stwo}{S^2}
\newcommand{\Sone}{S^1}
\newcommand{\SU}{\mathrm{SU}(2)}
\newcommand{\SO}{\mathrm{SO}}
\newcommand{\R}{\mathbb{R}}
\newcommand{\C}{\mathbb{C}}
\newcommand{\Z}{\mathbb{Z}}
\newcommand{\lap}{\Delta_{\Sthree}}
\newcommand{\bra}[1]{\langle #1 |}
\newcommand{\ket}[1]{| #1 \rangle}
\newcommand{\braket}[2]{\langle #1 | #2 \rangle}
\newcommand{\norm}[1]{\left\| #1 \right\|}

% ─────────────────────────────────────────────────────────────
\title{Hydrogen Orbital Structure from Harmonic Analysis on $S^3$}
\author{}
\date{}

\begin{document}
\maketitle

% ── Abstract ─────────────────────────────────────────────────
\begin{abstract}
We show that every hydrogen orbital shape emerges from harmonic analysis
on the 3-sphere~$\Sthree$. An ergodic trajectory on~$\Sthree$ (generated
here by composing three $\SU$ rotations with $\Z_3$ cyclic symmetry) is
weighted by eigenfunctions of the Laplacian~$\lap$. Each eigenfunction
selects a different orbital: $s$, $p$, $d$, and~$f$, for all principal
quantum numbers~$n$. No Schr\"odinger equation or Coulomb potential is used.
Three results go beyond Fock's 1935 momentum-space
correspondence~\cite{fock1935}: (i)~the shell capacity $2n^2$ follows
directly from the eigenspace dimension $(k{+}1)^2$ at Laplacian degree
$k = n{-}1$, without invoking the Coulomb potential; (ii)~the Aufbau energy
ordering follows from the strictly increasing eigenvalue sequence
$-k(k{+}2)$: $0, 3, 8, 15, \ldots$, making the filling order an inevitable
consequence of~$\Sthree$ geometry; (iii)~the radial wavefunctions are
Gegenbauer polynomials $C_k^{l+1}(\cos\chi)$ on the compact
interval~$[0,\pi]$, rather than Laguerre polynomials on~$[0,\infty)$ ---
both give $n{-}l{-}1$ radial nodes, but the node positions differ due to
$\Sthree$ compactification, yielding a falsifiable prediction. An interactive
WebGPU demonstration accompanies this paper. The orbital shape is not
imposed~--- it is inevitable geometry.
\end{abstract}

% ═════════════════════════════════════════════════════════════
\section{Introduction}
\label{sec:intro}
% ═════════════════════════════════════════════════════════════

The hydrogen atom has been the proving ground of quantum mechanics since
Bohr's original model~\cite{bohr1913}. The standard derivation begins with
the Schr\"odinger equation for an electron in a Coulomb
potential~\cite{schrodinger1926,griffiths2018}:
\begin{equation}
  -\frac{\hbar^2}{2m_e}\nabla^2\psi - \frac{e^2}{4\pi\epsilon_0 r}\psi
  = E\psi.
  \label{eq:schrodinger}
\end{equation}
Separation of variables in spherical coordinates yields the well-known
wavefunctions $\psi_{nlm}(r,\theta,\phi)$, characterized by quantum
numbers~$n$, $l$, and~$m$. This is a PDE boundary value problem: the
potential determines the equation, separation of variables determines the
coordinate system, and regularity conditions select the allowed solutions.

We present an alternative route to the \emph{same} orbital shapes that uses
none of this machinery. Our construction requires only:
\begin{enumerate}[nosep]
  \item The 3-sphere $\Sthree$ viewed as the Lie group $\SU$,
  \item Three independent $\SU$ rotations (quarks) with $\Z_3$ cyclic symmetry,
  \item Eigenfunctions of the Laplacian $\lap$ as spectral filters.
\end{enumerate}

The key idea is this: the electron's path through space is determined by the
composed rotation of three quarks. This path is ergodic on~$\Sthree$ ---
it eventually visits every region. Each orbital corresponds to a different
\emph{harmonic weight function} evaluated at the electron's $\Sthree$
position. The weight determines which regions of the trajectory contribute
to the visible probability cloud. Quantum numbers are not imposed by hand;
they emerge as eigenvalue labels of~$\lap$.

This is not merely an analogy. The connection between hydrogen and~$\Sthree$
has deep mathematical roots. Fock~\cite{fock1935} showed in 1935 that the
hydrogen atom possesses an $\SO(4)$ symmetry, and that momentum-space
wavefunctions live naturally on~$\Sthree$. Pauli~\cite{pauli1926} had
already exploited the $\SO(4)$ symmetry algebraically. What we demonstrate
is that this $\Sthree$ structure is not merely a mathematical convenience
but the \emph{geometric origin} of orbital shapes, accessible through a
simple quark rotation model and spectral filtering.

The paper is organized as follows. \Cref{sec:hopf} reviews the Hopf
fibration and $\Sthree$ as $\SU$. \Cref{sec:quarks} introduces the quark
rotation model and proves ergodicity of the electron trajectory.
\Cref{sec:harmonics} --- the core section --- develops $\Sthree$ harmonics
as spectral filters, derives shell capacity and the Aufbau principle from
the Laplacian spectrum, and situates the work relative to Fock's 1935
correspondence. \Cref{sec:projection} describes stereographic projection.
\Cref{sec:radial} presents the Gegenbauer vs.\ Laguerre radial comparison
--- the falsifiable prediction. \Cref{sec:results} gives quantitative
results, and \Cref{sec:discussion} discusses limitations and the
relationship to standard quantum mechanics.


% ═════════════════════════════════════════════════════════════
\section{The Hopf Fibration}
\label{sec:hopf}
% ═════════════════════════════════════════════════════════════

\subsection{$\Sthree$ as $\SU$}

The 3-sphere $\Sthree = \{(y_0, y_1, y_2, y_3) \in \R^4 : y_0^2 + y_1^2 +
y_2^2 + y_3^2 = 1\}$ carries the structure of the Lie group $\SU$. Every
point can be written as a unit quaternion
\begin{equation}
  q = a + b\mathbf{i} + c\mathbf{j} + d\mathbf{k},
  \quad |q|^2 = a^2 + b^2 + c^2 + d^2 = 1,
  \label{eq:quaternion}
\end{equation}
or equivalently as a $2\times 2$ unitary matrix with determinant~1:
\begin{equation}
  q = \begin{pmatrix} \alpha & -\bar{\beta} \\ \beta & \bar{\alpha}
  \end{pmatrix}, \quad |\alpha|^2 + |\beta|^2 = 1,
  \label{eq:su2matrix}
\end{equation}
where $\alpha = a + bi$ and $\beta = c + di$. Group multiplication is
quaternion multiplication, which is associative but non-commutative.

\subsection{The Fibration $\Sone \hookrightarrow \Sthree \to \Stwo$}

The Hopf fibration~\cite{hopf1931} is the unique non-trivial $\Sone$-bundle
over~$\Stwo$. It maps each point of $\Sthree$ to a point on $\Stwo$ (the
\emph{base space}), and the preimage of each base point is a great circle
$\Sone$ (the \emph{fiber}).

Explicitly, a point on the Hopf fiber at base coordinates $(\theta, \phi)$
on $\Stwo$ and fiber position $\psi \in [0, 2\pi)$ is:
\begin{equation}
  p(\theta, \phi, \psi) = \bigl(
    \cos\tfrac{\theta}{2}\cos\psi,\;
    \cos\tfrac{\theta}{2}\sin\psi,\;
    \sin\tfrac{\theta}{2}\cos(\psi+\phi),\;
    \sin\tfrac{\theta}{2}\sin(\psi+\phi)
  \bigr).
  \label{eq:hopf_fiber}
\end{equation}
The parameter $\psi$ traces the $\Sone$ fiber, while $(\theta, \phi)$
identifies which fiber we are on. Every point of $\Sthree$ belongs to
exactly one fiber.

\subsection{$\SU$ Action on $\Sthree$}

Left multiplication by $q \in \SU$ defines an isometry of $\Sthree$:
\begin{equation}
  L_q : p \mapsto q \cdot p.
  \label{eq:left_action}
\end{equation}
This action is \emph{transitive} --- any two points of $\Sthree$ are related
by some rotation --- and \emph{free} --- only $q = 1$ fixes every point.
Crucially, left $\SU$ action preserves the Hopf fibration structure:
it maps fibers to fibers~\cite{urbantke2003,nakahara2003}.

The composition of two $\SU$ elements $q_1, q_2$ is given by quaternion
multiplication:
\begin{equation}
  (q_1 \cdot q_2)_k = \sum_{i,j} \epsilon_{kij}\, q_{1,i}\, q_{2,j},
  \label{eq:su2_mul}
\end{equation}
where the coefficients follow from the quaternion algebra
$\mathbf{i}^2 = \mathbf{j}^2 = \mathbf{k}^2 = \mathbf{ijk} = -1$.
In components $(a, b, c, d)$:
\begin{equation}
  q_1 \cdot q_2 = \begin{pmatrix}
    a_1 a_2 - b_1 b_2 - c_1 c_2 - d_1 d_2 \\
    a_1 b_2 + b_1 a_2 + c_1 d_2 - d_1 c_2 \\
    a_1 c_2 - b_1 d_2 + c_1 a_2 + d_1 b_2 \\
    a_1 d_2 + b_1 c_2 - c_1 b_2 + d_1 a_2
  \end{pmatrix}^{\!\!\top}\!.
  \label{eq:quat_mul}
\end{equation}


% ═════════════════════════════════════════════════════════════
\section{The Quark Rotation Model}
\label{sec:quarks}
% ═════════════════════════════════════════════════════════════

\subsection{Three $\SU$ Paths with $\Z_3$ Symmetry}

We define three time-dependent $\SU$ elements representing the three color
charges of a proton's quarks. Each quark follows a 4-frequency path on
$\Sthree$, and the three quarks are related by \emph{cyclic permutation}
of frequencies --- a $\Z_3$ symmetry.

Given four frequencies $\omega_1, \omega_2, \omega_3, \omega_4 > 0$, define
the \emph{4-frequency $\SU$ path}:
\begin{equation}
  \hat{q}(t; \omega_1, \omega_2, \omega_3, \omega_4) =
  \frac{1}{N(t)}\bigl(\cos\omega_1 t,\; \sin\omega_2 t,\;
  \cos\omega_3 t,\; \sin\omega_4 t\bigr),
  \label{eq:su2_4freq}
\end{equation}
where $N(t) = \sqrt{\cos^2\omega_1 t + \sin^2\omega_2 t + \cos^2\omega_3 t
+ \sin^2\omega_4 t}$ normalizes to $\Sthree$.

The three quark rotations are:
\begin{align}
  q_R(t) &= \hat{q}(t;\; \omega_1, \omega_2, \omega_3, \omega_4),
  \label{eq:quark_R} \\
  q_G(t) &= \hat{q}(t;\; \omega_2, \omega_3, \omega_4, \omega_1),
  \label{eq:quark_G} \\
  q_B(t) &= \hat{q}(-t;\; \omega_3, \omega_4, \omega_1, \omega_2).
  \label{eq:quark_B}
\end{align}
The subscripts $R$, $G$, $B$ denote the three color charges. The key
features are:
\begin{itemize}[nosep]
  \item \textbf{Cyclic permutation}: $q_G$ uses the same frequencies as
    $q_R$ but shifted cyclically by one position.
  \item \textbf{Counter-rotation}: $q_B$ evolves with \emph{reversed} time
    and a two-position shift, reflecting the anti-quark structure of the
    proton's third valence quark.
  \item \textbf{$\Z_3$ symmetry}: the three quarks are related by the cyclic
    group $\Z_3 = \{e, \omega, \omega^2\}$ acting on the frequency labels.
\end{itemize}

\subsection{The Electron Rotation}

The electron's trajectory is the \emph{composed} rotation:
\begin{equation}
  q(t) = q_R(t) \cdot q_G(t) \cdot q_B(t),
  \label{eq:electron}
\end{equation}
where $\cdot$ denotes $\SU$ group multiplication
(\cref{eq:quat_mul}). The electron does not have its own independent
dynamics --- its state is entirely determined by the three quark rotations.

The electron's position on $\Sthree$ at time $t$ is obtained by acting with
$q(t)$ on a reference fiber point:
\begin{equation}
  y(t) = q(t) \cdot p_0,
  \label{eq:electron_pos}
\end{equation}
where $p_0 = p(\theta_0, 0, \psi(t))$ is a point on the Hopf fiber at some
fixed base angle $\theta_0$, with fiber phase
$\psi(t) = \pi\varphi\, t$ and $\varphi = (\sqrt{5}-1)/2$ the golden ratio.

\begin{remark}
The golden ratio ensures that the fiber phase $\psi(t)$ is
\emph{incommensurate} with all four frequencies $\omega_i$. This is critical
for ergodic coverage of the full $\Sthree$, as we show next.
\end{remark}

\subsection{Ergodicity}

\begin{theorem}[Informal]
\label{thm:ergodic}
If $\omega_1, \omega_2, \omega_3, \omega_4$ are mutually incommensurate
(no rational linear relation $\sum n_i \omega_i = 0$ with
$n_i \in \Z$ not all zero), then the trajectory
$\{y(t) : t \geq 0\}$ is dense in $\Sthree$.
\end{theorem}

\begin{proof}[Sketch]
Each $\hat{q}(t; \boldsymbol{\omega})$ traces a path on $\Sthree$ whose
closure is a torus of dimension equal to the number of independent
frequencies. With four mutually incommensurate frequencies, each quark path
is dense in a 4-torus embedded in $\Sthree$ (after normalization, this
projects to a dense subset of $\Sthree$). The composed rotation
$q = q_R \cdot q_G \cdot q_B$ inherits this property: the product of
ergodic $\SU$-valued paths with independent frequency spectra is itself
ergodic on $\SU \cong \Sthree$. Adding the incommensurate fiber phase
$\psi(t)$ ensures coverage of the $\Sone$ fiber direction as well.
\end{proof}

The physical consequence: the electron \emph{eventually visits every region
of $\Sthree$}. Its time-averaged position distribution converges to the
uniform (Haar) measure on $\Sthree$. The orbital shape is then determined
entirely by the harmonic weight function --- the topic of the next section.


% ═════════════════════════════════════════════════════════════
\section{$\Sthree$ Harmonics and Orbital Shapes}
\label{sec:harmonics}
% ═════════════════════════════════════════════════════════════

This is the core of the paper. We show that the eigenfunctions of the
Laplacian on $\Sthree$, when used as weight functions on the ergodic
electron trajectory, reproduce every hydrogen orbital shape.

\subsection{The Laplacian on $\Sthree$ and Its Spectrum}

The Laplace--Beltrami operator $\lap$ on the unit $\Sthree$ with the round
metric has eigenvalues~\cite{bander1966,frankel2011}:
\begin{equation}
  \lap Y_k = -k(k+2)\, Y_k, \quad k = 0, 1, 2, \ldots
  \label{eq:eigenvalues}
\end{equation}
The degeneracy of the $k$-th eigenspace is $(k+1)^2$. Setting $k = n-1$
(the principal quantum number minus one):
\begin{equation}
  \text{dim}(\text{eigenspace at } k = n-1) = n^2.
  \label{eq:degeneracy}
\end{equation}
Including the spin-$\frac{1}{2}$ degree of freedom (two orientations of the
$\Sone$ fiber), the total capacity of shell $n$ is:
\begin{equation}
  \boxed{C_n = 2n^2.}
  \label{eq:shell_capacity}
\end{equation}
This is precisely the hydrogen shell capacity, derived here from the
spectrum of $\lap$ without any reference to the Coulomb potential or angular
momentum algebra.

The $n^2$-dimensional eigenspace decomposes under the $\SO(3)$ subgroup
(rotations of the base $\Stwo$) into subshells labeled by angular momentum
$l = 0, 1, \ldots, n-1$, each with dimension $2l+1$:
\begin{equation}
  n^2 = \sum_{l=0}^{n-1}(2l+1), \qquad
  \text{subshell capacity} = 2(2l+1).
  \label{eq:subshell}
\end{equation}

\begin{remark}
In the standard treatment, the shell capacity $2n^2$ follows from solving
the radial Schr\"odinger equation and counting angular momentum states. Here
it follows from a single fact: the Laplacian on $\Sthree$ has eigenspace
dimension $(k+1)^2$. The quantum numbers are \emph{labels}, not inputs.
\end{remark}

\subsection{Weight Functions as Spectral Filters}

Each orbital $(n, l, m)$ corresponds to a harmonic $Y_{nlm}$ on $\Sthree$.
We use its squared modulus $|Y_{nlm}(y)|^2$ as a \emph{weight function}
that determines how much each point of the electron trajectory contributes
to the probability cloud.

The essential construction:
\begin{enumerate}[nosep]
  \item The electron traces the \emph{same} ergodic path on $\Sthree$ for
    all orbitals.
  \item At each time step, the electron is at some point $y(t) \in \Sthree$.
  \item The weight $w(t) = |Y_{nlm}(y(t))|^2$ determines the opacity/density
    of the trail at that point.
  \item Where $w$ is large, the cloud is dense; where $w \approx 0$ (a nodal
    region), the cloud vanishes.
\end{enumerate}

\begin{proposition}
The time-averaged weighted density converges to $|Y_{nlm}|^2$ integrated
against the Haar measure on $\Sthree$.
\end{proposition}

\begin{proof}
By \cref{thm:ergodic}, the electron trajectory is equidistributed on
$\Sthree$ with respect to Haar measure $\mu$. For any continuous
function $f$ on $\Sthree$:
\[
  \lim_{T\to\infty} \frac{1}{T}\int_0^T f(y(t))\,dt
  = \int_{\Sthree} f\,d\mu.
\]
Taking $f = |Y_{nlm}|^2 \cdot \chi_U$ for any measurable set $U$ gives
the result.
\end{proof}

\subsection{The Complete Harmonic Table}

The following table gives the explicit weight functions for all orbitals
from $n=1$ through $n=4$, expressed in terms of the $\Sthree$ coordinates
$y = (y_0, y_1, y_2, y_3)$.

We use the notation $\rho^2 \coloneqq y_0^2 + y_1^2$ for the ``fiber
plane'' coordinates and note that the Gegenbauer polynomials
$C_k^1(y_3)$ provide the radial structure.

\bigskip
\begin{table}[h!]
\centering
\caption{$\Sthree$ harmonic weight functions for hydrogen orbitals,
  $n = 1$ through $n = 4$. Each weight is $|Y_{nlm}(y)|^2$.}
\label{tab:harmonics}
\renewcommand{\arraystretch}{1.3}
\begin{tabular}{@{}cccp{7.5cm}@{}}
\toprule
$n$ & $l$ & $|m|$ & Weight function $w(y)$ \\
\midrule
1 & 0 & 0 & $1$ \\
\midrule
2 & 0 & 0 & $y_3^2$ \\
2 & 1 & 0 & $y_2^2$ \\
2 & 1 & 1 & $\rho^2 = y_0^2 + y_1^2$ \\
\midrule
3 & 0 & 0 & $(4y_3^2 - 1)^2$
  \quad {\footnotesize $= [C_2^1(y_3)]^2$} \\
3 & 1 & 0 & $y_3^2 \, y_2^2$ \\
3 & 1 & 1 & $y_3^2 \, \rho^2$ \\
3 & 2 & 0 & $(2y_2^2 - \rho^2)^2$ \\
3 & 2 & 1 & $y_2^2 \, \rho^2$ \\
3 & 2 & 2 & $\rho^4$ \\
\midrule
4 & 0 & 0 & $(8y_3^3 - 4y_3)^2$
  \quad {\footnotesize $= [C_3^1(y_3)]^2$} \\
4 & 1 & 0 & $(12y_3^2 - 2)^2 \, y_2^2$
  \quad {\footnotesize $[C_2^2(y_3)]^2 \, y_2^2$} \\
4 & 1 & 1 & $(12y_3^2 - 2)^2 \, \rho^2$ \\
4 & 2 & 0 & $y_3^2\,(2y_2^2 - \rho^2)^2$ \\
4 & 2 & 1 & $y_3^2 \, y_2^2 \, \rho^2$ \\
4 & 2 & 2 & $y_3^2 \, \rho^4$ \\
4 & 3 & 0 & $y_2^2\,(2y_2^2 - 3\rho^2)^2$ \\
4 & 3 & 1 & $(4y_2^2 - \rho^2)^2\,\rho^2$ \\
4 & 3 & 2 & $y_2^2 \, \rho^4$ \\
4 & 3 & 3 & $\rho^6$ \\
\bottomrule
\end{tabular}
\end{table}

\subsection{Structure of the Weight Functions}

Several patterns are visible in \cref{tab:harmonics}:

\paragraph{Radial structure from Gegenbauer polynomials.}
The $s$-orbital ($l=0$) weights for each shell $n$ are squared Gegenbauer
polynomials $[C_{n-1}^1(y_3)]^2$. These polynomials have $n-1$ zeros in
$y_3 \in (-1, 1)$, creating $n-1$ nodal surfaces that project to concentric
nodal shells in $\R^3$. For example:
\begin{itemize}[nosep]
  \item $n=1$: $w = 1$ (no nodes, spherical),
  \item $n=2$: $w = y_3^2$ (one node at $y_3 = 0$),
  \item $n=3$: $w = (4y_3^2 - 1)^2$ (nodes at $y_3 = \pm\tfrac{1}{2}$),
  \item $n=4$: $w = (8y_3^3 - 4y_3)^2$ (nodes at $y_3 = 0, \pm\tfrac{1}{\sqrt{2}}$).
\end{itemize}

\paragraph{Angular structure from the fiber plane.}
The coordinates $y_0, y_1$ parameterize the Hopf fiber direction, while
$y_2$ is the ``polar'' direction on the base $\Stwo$. Increasing $|m|$
corresponds to increasing powers of $\rho^2 = y_0^2 + y_1^2$, which
selects regions near the equator of the fibration. The maximum angular
momentum state $|m| = l$ has weight $\propto \rho^{2l}$, concentrated
entirely on the fiber plane.

\paragraph{Factorized structure.}
Most weights factor as (radial part) $\times$ (angular part), reflecting the
separation of the $\Sthree$ Laplacian into a ``radial'' piece
(Gegenbauer in $y_3$) and an angular piece (spherical harmonic in
$y_0, y_1, y_2$). This factorization is the $\Sthree$ analog of the
separation $\psi_{nlm} = R_{nl}(r) \cdot Y_l^m(\theta, \phi)$ in the
standard treatment.

\subsection{Shell Capacity and the Aufbau Principle}
\label{sec:aufbau}

The eigenvalue spectrum of~$\lap$ determines not only the orbital shapes but
also how electrons fill them. The $k$-th eigenvalue $-k(k{+}2)$ is
\emph{strictly increasing} in~$k$, so there is a unique energy ordering of
the harmonic degrees. Setting $k = n{-}1$:

\begin{center}
\renewcommand{\arraystretch}{1.3}
\begin{tabular}{ccccl}
\toprule
$k$ & Eigenvalue $-k(k{+}2)$ & Degeneracy $(k{+}1)^2$
  & Capacity $2(k{+}1)^2$ & Subshell decomposition \\
\midrule
0 & $0$   & 1  & 2  & $1s^2$ \\
1 & $-3$  & 4  & 8  & $2s^2 + 2p^6$ \\
2 & $-8$  & 9  & 18 & $3s^2 + 3p^6 + 3d^{10}$ \\
3 & $-15$ & 16 & 32 & $4s^2 + 4p^6 + 4d^{10} + 4f^{14}$ \\
\bottomrule
\end{tabular}
\end{center}

The Aufbau principle --- electrons fill shells in order of increasing
energy --- is typically presented as an empirical rule justified by the
Madelung $(n{+}l)$ ordering. Here it is a theorem: the eigenvalue sequence
$0, 3, 8, 15, \ldots$ is strictly increasing, so there is exactly one way
to fill from the lowest available state upward. Within each shell, the
$(k{+}1)^2$ degeneracy decomposes under $\SO(3)$ into subshells
$l = 0, 1, \ldots, k$, each with multiplicity $2l{+}1$. The factor of~2
from spin (two orientations of the $\Sone$ fiber) yields the familiar
capacities: $2, 8, 18, 32$.

\begin{remark}
For hydrogen (one electron), all subshells within a shell are degenerate ---
a consequence of the full $\SO(4)$ symmetry. For multi-electron atoms,
inter-electron repulsion lifts this degeneracy and the $(n{+}l)$ Madelung
ordering emerges. The spectral theory gives the \emph{shell} structure;
many-body effects determine the \emph{subshell} ordering within a shell.
\end{remark}

\subsection{Coulomb Orthogonality and State Exclusion}
\label{sec:coulomb_ortho}

Why do electrons prefer to occupy \emph{different} harmonics rather than
crowd into the same one? The answer lies in orthogonality of the $\Sthree$
eigenfunctions combined with Coulomb repulsion.

Consider two electrons, each weighted by an $\Sthree$ harmonic. The
electrostatic repulsion between their projected charge distributions is:
\begin{equation}
  E_{12} = \int\!\!\int \frac{\rho_1(\mathbf{r}_1)\,\rho_2(\mathbf{r}_2)}
  {|\mathbf{r}_1 - \mathbf{r}_2|}\,d^3\!r_1\,d^3\!r_2,
  \label{eq:coulomb_repulsion}
\end{equation}
where $\rho_i \propto |Y_{n_i l_i m_i}|^2$ is the projected density from
the $i$-th harmonic. When both electrons occupy the \emph{same} harmonic,
$\rho_1 = \rho_2$, the overlap is maximal and $E_{12}$ is largest. When they
occupy \emph{different} harmonics, the orthogonality of the eigenfunctions
on $\Sthree$ reduces the spatial overlap: numerical computation shows that
cross-harmonic repulsion is roughly half the same-harmonic
value~(see~\cref{sec:results}).

This provides a geometric explanation for why electrons fill orthogonal
states: it is energetically favorable. However, we note carefully that this
is \emph{consistent with} the Pauli exclusion principle --- it is not a
derivation of fermion statistics from geometry. The exclusion principle is an
independent postulate about the antisymmetry of multi-particle wavefunctions;
our result shows that Coulomb repulsion on $\Sthree$ provides an additional,
classical energetic incentive toward the same outcome.

\subsection{The Fock Correspondence: Foundation and Extensions}
\label{sec:fock}

The connection between $\Sthree$ harmonics and hydrogen wavefunctions is not
a new observation. Fock~\cite{fock1935} proved in 1935 that the Fourier
transform of the hydrogen wavefunction $\psi_{nlm}$, after stereographic
projection, \emph{is} a harmonic on $\Sthree$:
\begin{equation}
  \tilde{\psi}_{nlm}(\mathbf{p})
  \xrightarrow{\text{stereo}} Y_{nlm}(y) \in L^2(\Sthree).
  \label{eq:fock}
\end{equation}
Mathematically, this arises because the Coulomb problem in momentum space is
equivalent to the Laplacian eigenvalue problem on $\Sthree$. The hydrogen
atom's ``hidden'' $\SO(4)$ symmetry~\cite{pauli1926,bander1966} is
precisely the isometry group of $\Sthree$. This is an exact mathematical
result, and the orbital shapes we generate are a direct consequence of it.

Our contribution is \emph{not} the discovery of this equivalence. What is
new in this work falls into three categories:
\begin{enumerate}[nosep]
  \item \textbf{Ergodic sampling model}: the quark rotation construction
    (\cref{sec:quarks}) provides a concrete mechanism that generates dense
    $\Sthree$ trajectories. The frequency-independence result
    (\cref{sec:results}) proves the orbital shapes are geometric invariants
    of $\Sthree$, not artifacts of the particular parameterization --- any
    ergodic flow on $\Sthree$ produces the same result.
  \item \textbf{Shell capacity and Aufbau from spectral theory}: the shell
    capacity $2n^2$ and the energy ordering of subshells follow from the
    eigenspace structure of~$\lap$ alone, without solving the Schr\"odinger
    equation or invoking the Coulomb potential (\cref{sec:aufbau}).
  \item \textbf{Gegenbauer vs.\ Laguerre prediction}: the radial
    wavefunctions on compact $\Sthree$ (Gegenbauer polynomials) differ from
    those on flat $\R^3$ (Laguerre polynomials) in node positions, yielding
    a falsifiable compactification signature (\cref{sec:radial}).
\end{enumerate}

\begin{remark}
The quark rotation is one of infinitely many ways to generate an ergodic
flow on $\Sthree$. Three independent irrational rotations, a geodesic flow
with incommensurate periods, or any topologically mixing map on $\SU$ would
produce the same orbital shapes after harmonic weighting. The $\Z_3$
structure is a \emph{physically motivated} choice, not a mathematical
necessity.
\end{remark}


% ═════════════════════════════════════════════════════════════
\section{Stereographic Projection and the Probability Cloud}
\label{sec:projection}
% ═════════════════════════════════════════════════════════════

\subsection{Projection from $\Sthree$ to $\R^3$}

We project $\Sthree$ points to $\R^3$ via stereographic projection from
the south pole $(0, 0, 0, -1)$:
\begin{equation}
  \pi_S : (y_0, y_1, y_2, y_3) \mapsto
  \left(\frac{y_0}{1+y_3},\; \frac{y_1}{1+y_3},\; \frac{y_2}{1+y_3}\right).
  \label{eq:stereo}
\end{equation}
This is conformal (angle-preserving) and maps the ``north pole''
$(0,0,0,1)$ to the origin. Points near the south pole map to large radii
in $\R^3$.

\subsection{Conformal Compactification}

For visualization, we apply a conformal compactification to keep the cloud
within a bounded region. Given a projected point at distance $r$ from the
origin:
\begin{equation}
  r_{\text{compact}} = R \cdot \frac{\tanh(r/R)}{r/R} \cdot r
  = R \tanh(r/R),
  \label{eq:compact}
\end{equation}
where $R$ is a compactification radius. This smoothly maps
$[0, \infty) \to [0, R)$, preserving the angular structure while bounding
the radial extent. The hyperbolic tangent ensures smooth behavior at both
small and large $r$.

\subsection{Radial Density and Exponential Decay}

For the $1s$ orbital ($w = 1$, no angular weighting), the time-averaged
density after projection follows an approximately exponential decay:
\begin{equation}
  \rho(r) \propto e^{-2r/a_0},
  \label{eq:radial}
\end{equation}
where $a_0$ emerges as a characteristic length from the $\Sthree$ curvature.
This matches the standard hydrogen $1s$ radial wavefunction
$|\psi_{1s}|^2 \propto e^{-2r/a_0}$.

We quantify the fit quality using the coefficient of determination $R^2$
from linear regression on $\ln\rho(r)$ vs.\ $r$. For the 1s orbital, the
interactive demo typically achieves $R^2 > 0.98$ after $\sim 10^5$
trajectory samples.

Higher orbitals exhibit the expected nodal structure: $ns$ orbitals show
$n-1$ radial nodes (zeros of the Gegenbauer polynomial), while $np$, $nd$,
$nf$ orbitals show angular nodes that break spherical symmetry into the
characteristic lobe patterns.


% ═════════════════════════════════════════════════════════════
\section{Radial Structure: Gegenbauer vs.\ Laguerre}
\label{sec:radial}
% ═════════════════════════════════════════════════════════════

The orbital \emph{shapes} on $\Sthree$ match standard hydrogen by the Fock
correspondence. But the radial wavefunctions differ in a physically
meaningful way, because $\Sthree$ is compact and $\R^3$ is not. This
difference yields a falsifiable prediction.

\subsection{Two Radial Bases}

In the standard treatment, the radial wavefunction of hydrogen is an
associated Laguerre polynomial on the half-line~$[0,\infty)$:
\begin{equation}
  R_{nl}(r) \propto \left(\frac{2r}{na_0}\right)^{\!l}
  L_{n-l-1}^{2l+1}\!\left(\frac{2r}{na_0}\right)
  e^{-r/(na_0)}.
  \label{eq:laguerre_radial}
\end{equation}
On $\Sthree$, the radial coordinate is the polar angle
$\chi \in [0, \pi]$ (equivalently $y_3 = \cos\chi$), and the radial
wavefunction is a Gegenbauer polynomial on the compact interval~$[-1, 1]$:
\begin{equation}
  \mathcal{R}_{nl}(\chi) \propto
  C_{n-l-1}^{l+1}(\cos\chi)\,\sin^l\!\chi.
  \label{eq:gegenbauer_radial}
\end{equation}
Both are eigenfunctions of their respective radial operators, and both are
members of the same orthogonal polynomial family (Gegenbauer polynomials
generalize Legendre polynomials on compact domains; Laguerre polynomials
are the flat-space limit).

\subsection{Matching Node Counts, Different Node Positions}

Both radial wavefunctions have exactly $n{-}l{-}1$ nodes --- this is a
consequence of the Sturm--Liouville theory that applies to both the compact
and non-compact cases. The number of radial nodes is a topological invariant:
it counts zero-crossings of an eigenfunction of a second-order operator.

However, the node \emph{positions} differ. On $\Sthree$, the nodes of
$C_{n-l-1}^{l+1}(\cos\chi)$ are compressed toward the origin relative to
the Laguerre nodes, because the compact geometry wraps back at
$\chi = \pi$. Under stereographic projection
$r = \tan(\chi/2)$, this compactification maps the $\Sthree$ nodes to
finite radii, while the outermost Laguerre nodes extend to
$r \sim n^2 a_0$. For example, for the $3s$ orbital ($n{-}l{-}1 = 2$
nodes):

\begin{center}
\renewcommand{\arraystretch}{1.2}
\begin{tabular}{ccc}
\toprule
Node & $\Sthree$ position ($r/a_0$ via stereo) & Flat $\R^3$ position ($r/a_0$) \\
\midrule
1 & $\approx 0.55$ & $\approx 1.90$ \\
2 & $\approx 2.08$ & $\approx 7.10$ \\
\bottomrule
\end{tabular}
\end{center}

The $\Sthree$ nodes are systematically shifted inward. The outer node is
particularly compressed: flat-space hydrogen extends to large radii where
$\Sthree$ has already wrapped past $\chi = \pi/2$.

\subsection{A Falsifiable Prediction}

This constitutes a testable prediction. If physical space has $\Sthree$
topology at some curvature scale, then precision measurements of hydrogen
orbital cross-sections --- particularly the radial probability distribution
$P(r) = |R_{nl}(r)|^2 r^2$ at large~$r$ --- should show systematic inward
shifts of node positions relative to flat-space Laguerre predictions. The
effect is strongest for high-$n$ orbitals where the outermost lobes probe
the largest radii.

Current spectroscopic precision is insufficient to detect this
compactification signature unless the $\Sthree$ radius is comparable to
$n^2 a_0 \sim 10^{-10}$~m. For cosmological-scale $\Sthree$ curvature, the
shifts are unobservably small. The prediction is nonetheless well-defined
and distinguishes the $\Sthree$ framework from flat-space quantum mechanics
in principle.


% ═════════════════════════════════════════════════════════════
\section{Results}
\label{sec:results}
% ═════════════════════════════════════════════════════════════

\subsection{Orbital Shapes}

All 20 hydrogen orbitals from $n=1$ through $n=4$ were generated using the
interactive WebGPU demonstration. In every case, the probability cloud
matches the expected orbital shape:

\begin{itemize}[nosep]
  \item \textbf{$1s$}: Spherically symmetric. $R^2 > 0.98$ for exponential
    radial fit.
  \item \textbf{$2s$}: Spherical with one radial node.
  \item \textbf{$2p$ ($m=0$)}: Dumbbell shape aligned with the $y_2$ axis.
  \item \textbf{$2p$ ($|m|=1$)}: Torus (doughnut) in the fiber plane.
  \item \textbf{$3d$ ($m=0$)}: Cloverleaf pattern.
  \item \textbf{$3d$ ($|m|=2$)}: Torus with four-fold angular modulation.
  \item \textbf{$4f$ ($|m|=3$)}: Six-lobed structure with maximum angular
    momentum.
\end{itemize}

The key observation: \emph{all orbital shapes emerge from the same quark
rotation path}. The quark frequencies $\omega_1, \ldots, \omega_4$ affect
only the rate of convergence, not the final shape. Randomizing the
frequencies (while maintaining incommensurability) produces identical
orbital shapes --- the geometry is universal.

\subsection{Frequency Independence}

The interactive demo includes a ``Randomize'' function that generates new
incommensurate frequency sets. Repeated randomization confirms that:
\begin{enumerate}[nosep]
  \item The orbital shapes are independent of the specific frequency values.
  \item The convergence rate depends on how ``far from rational'' the
    frequencies are (more irrational $\Rightarrow$ faster coverage of
    $\Sthree$).
  \item The exponential fit parameters ($a_0$, $R^2$) converge to the same
    values regardless of frequency choice.
\end{enumerate}

This universality is a direct consequence of ergodicity: any set of
mutually incommensurate frequencies produces a dense trajectory on $\Sthree$,
and the time average converges to the Haar measure. The orbital shape is a
property of the \emph{harmonic filter}, not of the \emph{trajectory}.

\subsection{Quantitative Shell and Subshell Structure}

The following capacities are exact consequences of the $\Sthree$ Laplacian
spectrum:

\begin{center}
\renewcommand{\arraystretch}{1.2}
\begin{tabular}{cccl}
\toprule
Shell $n$ & Eigenvalue $k(k+2)$ & Degeneracy & Subshells \\
\midrule
1 & 0 & 1 & $1s$ \\
2 & 3 & 4 & $2s + 2p$ \\
3 & 8 & 9 & $3s + 3p + 3d$ \\
4 & 15 & 16 & $4s + 4p + 4d + 4f$ \\
\bottomrule
\end{tabular}
\end{center}

Including spin: $C_n = 2n^2$ gives 2, 8, 18, 32 for $n = 1, 2, 3, 4$ ---
exactly the hydrogen shell capacities.

\subsection{Quantitative Comparison with Standard Quantum Mechanics}

\Cref{tab:comparison} summarizes the quantitative agreement between the
$\Sthree$ construction and standard hydrogen quantum mechanics across all
observable categories. Numerical values are from the accompanying test
suite (\texttt{hydrogen\_comparison\_test.go}, \texttt{hopf\_states\_test.go}).

\begin{table}[h!]
\centering
\caption{Comparison of $\Sthree$ predictions with standard quantum mechanics
  for hydrogen orbitals, $n = 1$ through $n = 4$.}
\label{tab:comparison}
\renewcommand{\arraystretch}{1.3}
\begin{tabular}{@{}p{3.2cm}p{3.8cm}p{3.5cm}c@{}}
\toprule
Observable & $\Sthree$ prediction & Standard QM & Match? \\
\midrule
Shell capacity &
  $2(k{+}1)^2$ from $\lap$ degeneracy &
  $2n^2$ from Coulomb degeneracy &
  Exact \\
Energy ordering &
  $k(k{+}2)$: $0, 3, 8, 15$ &
  $-13.6/n^2$: same ordering &
  Same \\
Radial nodes &
  $n{-}l{-}1$ (Gegenbauer zeros) &
  $n{-}l{-}1$ (Laguerre zeros) &
  Exact \\
Node positions &
  Compressed (compact $\Sthree$) &
  Standard (flat $\R^3$) &
  \textbf{Differs} \\
Angular patterns &
  $\Sthree$ polynomials in $(y_0, y_1, y_2)$ &
  Spherical harmonics $Y_l^m$ &
  $r > 0.99$ \\
$1s$ radial fit &
  $R^2 > 0.98$ &
  Exact exponential &
  Matches \\
Coulomb repulsion &
  Same-harmonic ${\approx}\,2\times$ cross &
  Same (Hund's rules) &
  Consistent \\
\bottomrule
\end{tabular}
\end{table}

The single entry marked \textbf{Differs} --- radial node positions --- is
the falsifiable prediction discussed in \cref{sec:radial}. All other entries
match exactly or to numerical precision ($r > 0.99$ angular correlation
across all 20 states tested).


% ═════════════════════════════════════════════════════════════
\section{Discussion}
\label{sec:discussion}
% ═════════════════════════════════════════════════════════════

\subsection{Comparison with the Standard Approach}

The standard derivation~\cite{griffiths2018} of hydrogen orbitals requires:
\begin{enumerate}[nosep]
  \item The Schr\"odinger equation (\cref{eq:schrodinger}),
  \item The Coulomb potential $V(r) = -e^2/(4\pi\epsilon_0 r)$,
  \item Separation of variables in spherical coordinates,
  \item Regularity conditions at $r = 0$ and $r \to \infty$,
  \item The resulting radial equation solved by associated Laguerre
    polynomials.
\end{enumerate}

Our construction replaces all five steps with:
\begin{enumerate}[nosep]
  \item Three $\SU$ paths with $\Z_3$ symmetry (the quark model),
  \item Eigenfunctions of $\lap$ (the harmonic filter).
\end{enumerate}

The two approaches produce the same orbital shapes because they describe
the same mathematical object from different perspectives: the Schr\"odinger
approach works in position/momentum space where the $\SO(4)$ symmetry is
hidden; the $\Sthree$ approach works directly on the symmetry manifold
where the structure is manifest~\cite{fock1935,bander1966}.

\subsection{What Is New Beyond Fock}

The Fock correspondence~\cite{fock1935} is the mathematical foundation of
this work: it establishes that $\Sthree$ harmonics \emph{are} hydrogen
wavefunctions in momentum space. Our contributions build on this foundation
in four directions:

\begin{enumerate}[nosep]
  \item \textbf{Shell capacity from spectral theory}: the capacity $2n^2$
    follows from the eigenspace dimension $(k{+}1)^2$ of $\lap$ without
    solving the Schr\"odinger equation or analyzing the Coulomb
    potential. Fock's paper uses the correspondence to simplify the hydrogen
    \emph{calculation}; we use the Laplacian spectrum to \emph{derive} the
    shell structure as a theorem about $\Sthree$.
  \item \textbf{Aufbau from eigenvalue ordering}: the strictly increasing
    sequence $k(k{+}2)$ determines the energy ordering. This is a
    consequence of $\Sthree$ geometry, not an empirical rule.
  \item \textbf{Falsifiable radial prediction}: Gegenbauer radial
    wavefunctions on compact $\Sthree$ predict different node positions from
    Laguerre wavefunctions on flat $\R^3$, providing a testable signature of
    compactification (\cref{sec:radial}).
  \item \textbf{Universal orbital shapes from ergodic flows}: frequency
    independence demonstrates that \emph{any} ergodic $\Sthree$ trajectory
    produces the same orbitals. The quark rotation model is one such flow;
    the shapes are properties of the geometry, not the parameterization.
\end{enumerate}

\subsection{Limitations}

Several important limitations should be stated clearly:

\begin{itemize}[nosep]
  \item \textbf{Energy eigenvalues}: The $E_n = -13.6/n^2$~eV spectrum
    requires the Coulomb potential or its $\Sthree$ equivalent (the
    curvature of stereographic projection). Our construction produces
    the shapes but not the energies.
  \item \textbf{The quark rotation is an ansatz}: the $\Z_3$ quark model
    is one of many possible ergodic flows on $\Sthree$. Frequency
    independence (\cref{sec:results}) proves the shapes do not depend on
    this choice, but the physical motivation for \emph{why} the flow should
    be $\Sthree$-ergodic is not determined within this paper.
  \item \textbf{The fiber squashing parameter $\beta$}: the broader theory
    uses a parameter $\beta$ that controls the Berger sphere geometry. This
    paper operates at $\beta = 1$ (the round $\Sthree$); the effects of
    $\beta \neq 1$ on orbital shapes are not addressed here.
  \item \textbf{Multi-electron atoms}: Electron-electron repulsion breaks
    the $\SO(4)$ symmetry. The extension to many-electron systems requires
    a self-consistent field approach on $\Sthree$.
  \item \textbf{Relativistic corrections}: Fine structure, spin-orbit
    coupling, and the Lamb shift require extending the geometry to include
    time-dependent effects on the fibration.
  \item \textbf{Dynamics}: Our construction shows that the orbital shapes
    are \emph{kinematically} inevitable from $\Sthree$ geometry. Whether
    the \emph{dynamics} (time evolution, transition amplitudes) can also be
    derived from $\Sthree$ geometry alone remains an open question.
\end{itemize}


% ═════════════════════════════════════════════════════════════
\section{Conclusion}
\label{sec:conclusion}
% ═════════════════════════════════════════════════════════════

Every hydrogen orbital shape --- $s$, $p$, $d$, and $f$, for all principal
quantum numbers --- emerges from harmonic analysis on $\Sthree$. An ergodic
trajectory (here generated by three $\SU$ rotations with $\Z_3$ symmetry)
is weighted by eigenfunctions of the Laplacian~$\lap$, and the orbital
shape appears. No Schr\"odinger equation is solved; no Coulomb potential is
imposed; no quantum numbers are put in by hand.

Building on Fock's 1935 momentum-space correspondence, this work establishes
three results: (i)~the shell capacity $2n^2$ follows from the eigenspace
dimension of~$\lap$, (ii)~the Aufbau ordering follows from the strictly
increasing eigenvalue sequence $k(k{+}2)$, and (iii)~the radial
wavefunctions (Gegenbauer on compact $\Sthree$ vs.\ Laguerre on flat
$\R^3$) differ in node positions, yielding a falsifiable compactification
prediction. Frequency independence demonstrates that these results are
properties of $\Sthree$ geometry, not artifacts of the parameterization.

The accompanying interactive WebGPU demonstration makes this visible in real
time: the probability cloud builds as the electron traverses its ergodic
path, and the characteristic shapes of hydrogen orbitals appear --- not
because a formula commands it, but because the geometry demands it.


% ═════════════════════════════════════════════════════════════
\appendix
% ═════════════════════════════════════════════════════════════

\section{Complete Harmonic Weight Functions}
\label{app:harmonics}

For reference, we list all 20 weight functions for $n = 1$ through $n = 4$,
using $\rho^2 = y_0^2 + y_1^2$.

\bigskip

\noindent\textbf{$n = 1$ (1 orbital):}
\begin{align*}
  w_{1,0,0}(y) &= 1
\end{align*}

\noindent\textbf{$n = 2$ (4 orbitals):}
\begin{align*}
  w_{2,0,0}(y) &= y_3^2 \\
  w_{2,1,0}(y) &= y_2^2 \\
  w_{2,1,\pm1}(y) &= \rho^2
\end{align*}

\noindent\textbf{$n = 3$ (9 orbitals):}
\begin{align*}
  w_{3,0,0}(y) &= (4y_3^2 - 1)^2 \\
  w_{3,1,0}(y) &= y_3^2\, y_2^2 \\
  w_{3,1,\pm1}(y) &= y_3^2\, \rho^2 \\
  w_{3,2,0}(y) &= (2y_2^2 - \rho^2)^2 \\
  w_{3,2,\pm1}(y) &= y_2^2\, \rho^2 \\
  w_{3,2,\pm2}(y) &= \rho^4
\end{align*}

\noindent\textbf{$n = 4$ (16 orbitals):}
\begin{align*}
  w_{4,0,0}(y) &= (8y_3^3 - 4y_3)^2 \\
  w_{4,1,0}(y) &= (12y_3^2 - 2)^2\, y_2^2 \\
  w_{4,1,\pm1}(y) &= (12y_3^2 - 2)^2\, \rho^2 \\
  w_{4,2,0}(y) &= y_3^2\,(2y_2^2 - \rho^2)^2 \\
  w_{4,2,\pm1}(y) &= y_3^2\, y_2^2\, \rho^2 \\
  w_{4,2,\pm2}(y) &= y_3^2\, \rho^4 \\
  w_{4,3,0}(y) &= y_2^2\,(2y_2^2 - 3\rho^2)^2 \\
  w_{4,3,\pm1}(y) &= (4y_2^2 - \rho^2)^2\, \rho^2 \\
  w_{4,3,\pm2}(y) &= y_2^2\, \rho^4 \\
  w_{4,3,\pm3}(y) &= \rho^6
\end{align*}

Note: counting $\pm m$ pairs separately gives
$1 + 4 + 9 + 16 = 30$ orbitals (without spin). With spin-$\frac{1}{2}$,
$2 + 8 + 18 + 32 = 60$ states, matching the first four hydrogen shells.


\section{Interactive Demo: Implementation Details}
\label{app:demo}

The interactive demonstration is implemented as a client-side WebGPU
application with no server-side computation.

\paragraph{Architecture.}
The application consists of four modules:
\begin{itemize}[nosep]
  \item \texttt{physics.js} --- SU(2) math, quark rotations, $\Sthree$
    harmonics, density accumulation, and exponential fit statistics.
  \item \texttt{renderer.js} --- WebGPU rendering pipeline with four layers:
    Hopf fibers, electron sphere, Z$_3$ nucleus, and probability cloud.
  \item \texttt{camera.js} --- Orbit camera with smooth inertial damping.
  \item \texttt{main.js} --- Frame loop, UI bindings, and live telemetry.
\end{itemize}

\paragraph{Density accumulation.}
At each time step, the electron's $\Sthree$ position is computed via
\cref{eq:electron,eq:electron_pos}. The harmonic weight at this position
is evaluated using the functions from \cref{tab:harmonics}. Points with
weight below $10^{-8}$ are discarded (nodal regions). Surviving points are
stereographically projected (\cref{eq:stereo}), compactified
(\cref{eq:compact}), and added to a ring buffer of up to $10^6$ points.

\paragraph{Rendering.}
The probability cloud is rendered as a point cloud with per-point opacity
proportional to the harmonic weight. Four WGSL shaders handle the distinct
visual layers. The GPU accumulates the visual density directly from the
physics buffer.

\paragraph{Telemetry.}
The right panel displays real-time values of all intermediate quantities:
the three quark rotors $(q_R, q_G, q_B)$, the composed rotation $q$, the
$\Sthree$ position $y$, the $\R^3$ projection, the fiber phase $\psi$, and
the harmonic weight $|Y_N|^2$. This makes the entire computation chain
transparent to the viewer.

The demo is available at the accompanying website and requires a
WebGPU-capable browser (Chrome 113+ or Edge 113+).


% ── Bibliography ─────────────────────────────────────────────
\bibliographystyle{unsrt}
\bibliography{bibliography}

\end{document}
